\documentclass[12pt, a4paper]{article}

% --- Paquetes básicos ---
\usepackage[utf8]{inputenc}
\usepackage[spanish]{babel}
\usepackage[margin=2.5cm]{geometry}
\usepackage{graphicx}
\usepackage{amsmath, amssymb}
\usepackage{fancyhdr}

% --- Configuración del Título ---
\title{Bloque 2: Simulación de Carga de Trabajo y Monitorización}
\author{Juan Reche Aguilar}
\date{\today}

% --- Configuración de Encabezados ---
\pagestyle{fancy}
\fancyhf{}
\rhead{Ingeniería de Servidores}
\lhead{Juan Reche Aguilar}
\cfoot{\thepage}

\begin{document}

\maketitle

\tableofcontents
\newpage

\section{Ejercicio 1:}

\begin{figure}[htbp]
    \centering

    \includegraphics[width=0.6\textwidth]{../capturas1/ejercicio1.png}

    \caption{Descripción del gráfico o figura del ejercicio.}

    \label{fig:mi_grafico}
\end{figure}

Como se puede observar en la Figura \ref{fig:mi_grafico}, los resultados coinciden con lo calculado en la teoría.

\end{document}